% Generated draft: requires mathematical and source review.
\documentclass[11pt]{article}
\usepackage[margin=25mm]{geometry}
\usepackage{fontspec}
\setmainfont{Latin Modern Roman}
\usepackage{amsmath,amssymb,mathtools,unicode-math}
\setmathfont{Latin Modern Math}
\usepackage{enumitem,needspace,xurl,hyperref}
\setlength{\parindent}{0pt}
\setlength{\parskip}{5pt}
\setlength{\emergencystretch}{3em}
\newcommand{\N}{\mathbb{N}}
\newcommand{\Z}{\mathbb{Z}}
\newcommand{\Q}{\mathbb{Q}}
\newcommand{\R}{\mathbb{R}}
\newcommand{\C}{\mathbb{C}}
\newcommand{\F}{\mathbb{F}}
\newcommand{\A}{\mathbb{A}}
\newcommand{\PP}{\mathbb{P}}
\newcommand{\E}{\mathbb{E}}
\newcommand{\eps}{\varepsilon}
\newcommand{\dd}{\,\mathrm{d}}
\begin{document}
\section*{UnsolvedMath ID: 11000145}
\textbf{Source problem code:} catalogue entry
\subsection{Definitions and mathematical statement}
For a Coxeter matrix \(M=(m_{ij})\), the Artin group has generators \(s_i\) and relations equating the two alternating words of length \(m_{ij}\) in \(s_i,s_j\), for each finite \(m_{ij}\); \(m_{ij}=\infty\) imposes no relation. The source asks which Artin groups admit injective homomorphisms into the surface mapping class group \(M(S)\) that are non-geometric in its preceding sense. It explicitly restricts attention to all exponents \(m_{ij}\) finite, because examples with infinite exponents follow from the missing Lemma 2(iii). Finite exponents do not mean that the associated Coxeter group is finite.

The source's geometric notion is not included; in the usual Dehn-twist setting it refers to sending standard generators to prescribed individual twists, but its exact conditions and equivalence conventions must be checked. Dropping geometricity removes the source's small-Artin-group restriction. Neither the surface range nor whether \(S\) is fixed is explicit.

As an example, the source describes the type-\(B_n\) diagram as type \(A_n\) with the first edge relabelled \(4\), and a map to type \(A_n\) taking the first generator to the square of the corresponding generator and the others to corresponding generators. In twist notation it prints
\[
 T_{\alpha_1}^2T_{\alpha_2}T_{\alpha_1}^2T_{\alpha_2}
 =T_{\alpha_2}T_{\alpha_1}^2T_{\alpha_2}T_{\alpha_1}^2.
\]
It also reports similar maps of other finite-type Artin groups into \(E_8\), using products of powers of a small number of twists, informally called almost geometric. The exact scope, formulas and injectivity of these maps are not supplied; existence of a homomorphism is not promoted to an embedding. The author only says these are the almost-geometric examples known to them, not that they exhaust all such maps.
\subsection{Short English statement}
Classify non-geometric Artin-group embeddings into the intended mapping class groups, restricting to finite pairwise Artin exponents and distinguishing that restriction from finite Coxeter type.
\subsection{Sources}
\begin{itemize}
\item \emph{Problems on Mapping Class Groups and Related Topics}, edited by Farb, unnumbered Question, supplied PDF page 131. \url{https://www.math.uchicago.edu/~farb/papers/mcgbook.pdf}. Recorded access 2026-07-29; extraction pdf-font-structure. The preceding geometric definition and Lemma 2(iii) are absent; no individual author or additional reference is supplied.
\item Local UnsolvedMath registry, \texttt{data/unsolvedmath/problems.json}, permanent ID \(11000145\), exact source code \texttt{catalogue entry}. Its OPEN-TRIAGE summary mentions Artin/right-Artin work without an exact citation; that mention does not remove the printed finite-exponent restriction.
\end{itemize}
\end{document}
